\section{Introduction}
\label{sec:intro}

The discovery of conventional superconductivity above 200~K in
hydrogen-rich compounds under pressure has transformed the search for
high-temperature superconductors. Sulfur hydride H$_3$S in the
Im$\bar{3}$m structure exhibits $T_c = 203$~K at
155~GPa~\cite{Drozdov2015}, while the clathrate superhydride
LaH$_{10}$ (Fm$\bar{3}$m) reaches $T_c = 250$~K at
170~GPa~\cite{Drozdov2019,Somayazulu2019}. Yttrium hexahydride
YH$_6$ ($T_c = 224$~K at 166~GPa)~\cite{Troyan2021} and calcium
hexahydride CaH$_6$ ($T_c = 215$~K at 172~GPa)~\cite{Ma2022} further
confirm that phonon-mediated pairing in dense hydrogen sublattices can
produce critical temperatures rivaling the cuprates. However, the
megabar pressures required for thermodynamic stability --- above
150~GPa in every confirmed case with $T_c > 200$~K --- confine these
materials to diamond anvil cells and preclude practical applications.
Reducing the stabilization pressure while preserving strong
electron-phonon coupling remains the central challenge of the
field~\cite{FloresLivas2020,Boeri2022}.

The chemical pre-compression strategy proposed by
Ashcroft~\cite{Ashcroft2004} offers a route toward this goal.
By embedding hydrogen within a host lattice of heavier elements, the
non-hydrogen sublattice provides internal chemical pressure that
compresses the H sublattice without commensurate external load.
This concept underlies all confirmed hydride superconductors and has
been extended to ternary systems, where a third element tunes
stability and electronic structure. The experimental confirmation of
LaBeH$_8$ ($T_c = 110$~K at 80~GPa)~\cite{Song2023} demonstrated
that light-element scaffolds can halve the required pressure relative
to binary superhydrides, albeit at the cost of a significant $T_c$
reduction. More recently, ambient-pressure predictions for Mg-based
hydrides~\cite{Lucrezi2024} and boron-carbon
clathrates~\cite{Liang2024} have pushed the stabilization pressure
toward zero, but with $T_c$ values below 160~K. The steep
pressure--$T_c$ tradeoff observed across these families motivates a
search for structural motifs that can sustain strong hydrogen-derived
electron-phonon coupling at near-ambient pressures.

Ternary perovskite hydrides MXH$_3$ (M = alkali metal, X = group-13
metal) represent a promising low-pressure platform. Du~\textit{et~al.}
predicted that several MXH$_3$ compounds in the cubic Pm$\bar{3}$m
structure are dynamically stable below 10~GPa, with the
corner-sharing XH$_6$ octahedral framework providing chemical
pre-compression of the hydrogen sublattice~\cite{Du2024}. The
perovskite architecture is well suited for this purpose: the rigid
XH$_6$ cage maintains short H--H distances and high hydrogen-derived
density of states at the Fermi level, while the A-site cation (Cs, Rb,
K) stabilizes the structure at pressures accessible to multi-anvil
synthesis. Initial estimates of electron-phonon coupling in this family
suggested $\lambda \sim 2$--$3$ and $T_c$ potentially exceeding 150~K,
but a systematic treatment including anharmonic corrections ---
essential for any hydrogen-rich system~\cite{Errea2020} --- had not
been performed.

In this Rapid Communication, we report a first-principles screening
of MXH$_3$ perovskite hydrides using density functional theory,
density functional perturbation theory, isotropic Eliashberg theory,
and stochastic self-consistent harmonic approximation (SSCHA)
anharmonic corrections. From six candidate compounds, three cubic
perovskites pass both thermodynamic ($E_{\mathrm{hull}} < 50$~meV/atom)
and dynamic stability criteria at $P \leq 10$~GPa. The top-ranked
candidate, CsInH$_3$ (Pm$\bar{3}$m), achieves an anharmonic
$T_c = 214$~K at $P = 3$~GPa ($\mu^* = 0.13$), matching the
experimental $T_c$ of H$_3$S at a pressure $30\times$ lower. The
structure is quantum-stabilized at 3~GPa by hydrogen zero-point
motion, with all SSCHA phonon frequencies real. These results identify
cubic perovskite hydrides as a route to $>$200~K superconductivity at
pressures within reach of large-volume press techniques.

%% CITATIONS NEEDED (for gpd-bibliographer):
%% - Drozdov2015 — Drozdov et al., H3S 203 K, Nature 525, 73 (2015)
%% - Drozdov2019 — Drozdov et al., LaH10 250 K, Nature 569, 528 (2019)
%% - Somayazulu2019 — Somayazulu et al., LaH10 confirmation, PRL 122, 027001 (2019)
%% - Troyan2021 — Troyan et al., YH6 224 K, Adv. Mater. 33, 2006832 (2021)
%% - Ma2022 — Ma et al., CaH6 215 K, arXiv:2103.16282
%% - FloresLivas2020 — Flores-Livas et al., review, Phys. Rep. 856, 1 (2020)
%% - Boeri2022 — Boeri et al., theory of hydride SC, J. Phys. Condens. Matter 34, 183002 (2022)
%% - Ashcroft2004 — Ashcroft, chemical precompression, PRL 92, 187002 (2004)
%% - Song2023 — Song et al., LaBeH8 at 80 GPa, PRL 130, 266001 (2023)
%% - Lucrezi2024 — Lucrezi et al., Mg-based hydride ambient pressure, PRL 132, 166001 (2024)
%% - Liang2024 — Liang et al., B-C clathrate hydrides, PRB 109, 184517 (2024)
%% - Du2024 — Du et al., MXH3 perovskite hydrides, Adv. Sci. 11, 2408370 (2024)
%% - Errea2020 — Errea et al., quantum effects in LaH10, Nature 578, 66 (2020)
